\documentclass[12pt,a4paper]{article}
\usepackage[utf8]{inputenc}
\usepackage[T1]{fontenc}
\usepackage[spanish]{babel}
\usepackage{geometry}
\usepackage{longtable}
\usepackage{array}
\usepackage{booktabs}
\usepackage{graphicx}
\usepackage{float}
\usepackage{hyperref}
\usepackage{xcolor}
\geometry{margin=2.5cm}
\title{Informe académico — Encuesta normal}
\author{SADI}
\date{}
\begin{document}
\maketitle
\section*{Datos generales}
\begin{itemize}
\item Dataset: Alumnos
\item Fecha de generación: 2026-03-12 16:14
\item Tipo de análisis: Encuesta normal
\end{itemize}
\section{Introducción}
El presente informe expone un análisis descriptivo e interpretativo de los datos correspondientes a una encuesta normal procesada en SADI. La base analizada estuvo conformada por 10 registros válidos para el análisis general. Asimismo, se detectó automáticamente la variable de segmentación 'Genero', lo que permitió explorar diferencias y asociaciones entre grupos. El informe resume las principales características del dataset, los resultados descriptivos, las asociaciones estadísticas detectadas y una interpretación preliminar de los hallazgos.
\section{Perfil de la encuesta}
\begin{itemize}
\item Total de participantes: 10
\item Variables analizadas: —
\item Variables categóricas: —
\item Variables numéricas: —
\item Segmento principal: Genero
\end{itemize}
\section{Descripción general del dataset}
El análisis descriptivo comprendió 7 variables, de las cuales 6 fueron numéricas, 0 categóricas, 0 temporales y 1 de otro tipo. En general, no se identificaron proporciones críticas de datos faltantes en las variables resumidas. Este panorama descriptivo permite identificar la estructura general de la base y orientar la interpretación posterior de los resultados.
\section{Resumen descriptivo por variable}
\subsection{Edad (numeric)}
N total: 10 | N válido: 10 | N faltante: 0 (0.0\textbackslash\{\}\%)
Media: 23.5 | Mediana: 23.0 | Mínimo: 21.0 | Máximo: 27.0 | Desv. estándar: 1.9578900207451218 | Q1: 22.0 | Q3: 24.75
\subsection{Genero (categorical)}
N total: 10 | N válido: 10 | N faltante: 0 (0.0\textbackslash\{\}\%)
\begin{longtable}{p{8cm}p{3cm}p{3cm}}
\toprule
Categoría & Frecuencia & \% \\
\midrule
\endhead
Femenino & 6 & 60.0 \\
Masculino & 4 & 40.0 \\
\bottomrule
\end{longtable}
\subsection{P1\_Claridad\_contenido (numeric)}
N total: 10 | N válido: 10 | N faltante: 0 (0.0\textbackslash\{\}\%)
Media: 4.2 | Mediana: 4.0 | Mínimo: 3.0 | Máximo: 5.0 | Desv. estándar: 0.7888106377466155 | Q1: 4.0 | Q3: 5.0
\subsection{P2\_Dominio\_docente (numeric)}
N total: 10 | N válido: 10 | N faltante: 0 (0.0\textbackslash\{\}\%)
Media: 4.2 | Mediana: 4.0 | Mínimo: 3.0 | Máximo: 5.0 | Desv. estándar: 0.6324555320336759 | Q1: 4.0 | Q3: 4.75
\subsection{P3\_Material\_apoyo (numeric)}
N total: 10 | N válido: 10 | N faltante: 0 (0.0\textbackslash\{\}\%)
Media: 4.3 | Mediana: 4.0 | Mínimo: 3.0 | Máximo: 5.0 | Desv. estándar: 0.6749485577105528 | Q1: 4.0 | Q3: 5.0
\subsection{P4\_Participacion\_clase (numeric)}
N total: 10 | N válido: 10 | N faltante: 0 (0.0\textbackslash\{\}\%)
Media: 3.9 | Mediana: 4.0 | Mínimo: 3.0 | Máximo: 5.0 | Desv. estándar: 0.7378647873726218 | Q1: 3.25 | Q3: 4.0
\subsection{P5\_Satisfaccion\_general (numeric)}
N total: 10 | N válido: 10 | N faltante: 0 (0.0\textbackslash\{\}\%)
Media: 4.4 | Mediana: 4.5 | Mínimo: 3.0 | Máximo: 5.0 | Desv. estándar: 0.6992058987801011 | Q1: 4.0 | Q3: 5.0
\section{Resultados e interpretación académica}
Participaron 10 personas en la encuesta.
La variable seleccionada como segmento fue 'Genero'.

En términos inferenciales, no se observaron asociaciones estadísticamente significativas de magnitud relevante entre las combinaciones evaluadas, lo cual podría reflejar una relativa homogeneidad en las respuestas o ausencia de diferencias marcadas entre grupos.

Entre las variables con mayor relevancia analítica destacan: P4\_Participacion\_clase.

Participaron 10 personas en la encuesta.
La variable seleccionada como segmento fue 'Genero'.
\section{Variables con mayor relevancia analítica}
\begin{itemize}
\item P4\_Participacion\_clase
\end{itemize}
\section{Comparación entre grupos}
El análisis comparativo entre grupos permitió identificar diferencias en las medias de algunas variables numéricas según la variable de segmentación detectada. En la variable 'Edad', el grupo 'Masculino' presentó la media más alta (M=24.25), mientras que 'Femenino' registró la más baja (M=23.00), evidenciando una diferencia alta de 1.25 puntos. En la variable 'P4\_Participacion\_clase', el grupo 'Femenino' presentó la media más alta (M=4.33), mientras que 'Masculino' registró la más baja (M=3.25), evidenciando una diferencia alta de 1.08 puntos. En la variable 'P1\_Claridad\_contenido', el grupo 'Femenino' presentó la media más alta (M=4.50), mientras que 'Masculino' registró la más baja (M=3.75), evidenciando una diferencia moderada de 0.75 puntos. En la variable 'P5\_Satisfaccion\_general', el grupo 'Femenino' presentó la media más alta (M=4.67), mientras que 'Masculino' registró la más baja (M=4.00), evidenciando una diferencia moderada de 0.67 puntos. En la variable 'P3\_Material\_apoyo', el grupo 'Femenino' presentó la media más alta (M=4.50), mientras que 'Masculino' registró la más baja (M=4.00), evidenciando una diferencia moderada de 0.50 puntos. Estas diferencias deben interpretarse como una aproximación descriptiva inicial y pueden complementarse posteriormente con análisis inferenciales más específicos.
\section{Gráficos del análisis}
\begin{figure}[H]
\centering
\includegraphics[width=0.9\textwidth]{C:/sadi/static/plots/ds4_normal_Edad_hist.png}
\caption{ds4\_normal\_Edad\_hist.png}
\end{figure}
\begin{figure}[H]
\centering
\includegraphics[width=0.9\textwidth]{C:/sadi/static/plots/ds4_normal_Genero_bar.png}
\caption{ds4\_normal\_Genero\_bar.png}
\end{figure}
\begin{figure}[H]
\centering
\includegraphics[width=0.9\textwidth]{C:/sadi/static/plots/ds4_normal_P1_Claridad_contenido_hist.png}
\caption{ds4\_normal\_P1\_Claridad\_contenido\_hist.png}
\end{figure}
\begin{figure}[H]
\centering
\includegraphics[width=0.9\textwidth]{C:/sadi/static/plots/ds4_normal_P2_Dominio_docente_hist.png}
\caption{ds4\_normal\_P2\_Dominio\_docente\_hist.png}
\end{figure}
\begin{figure}[H]
\centering
\includegraphics[width=0.9\textwidth]{C:/sadi/static/plots/ds4_normal_P3_Material_apoyo_hist.png}
\caption{ds4\_normal\_P3\_Material\_apoyo\_hist.png}
\end{figure}
\begin{figure}[H]
\centering
\includegraphics[width=0.9\textwidth]{C:/sadi/static/plots/ds4_normal_P4_Participacion_clase_hist.png}
\caption{ds4\_normal\_P4\_Participacion\_clase\_hist.png}
\end{figure}
\begin{figure}[H]
\centering
\includegraphics[width=0.9\textwidth]{C:/sadi/static/plots/ds4_normal_P5_Satisfaccion_general_hist.png}
\caption{ds4\_normal\_P5\_Satisfaccion\_general\_hist.png}
\end{figure}
\begin{figure}[H]
\centering
\includegraphics[width=0.9\textwidth]{C:/sadi/static/plots/ds4_group_mean_genero_edad.png}
\caption{ds4\_group\_mean\_genero\_edad.png}
\end{figure}
\begin{figure}[H]
\centering
\includegraphics[width=0.9\textwidth]{C:/sadi/static/plots/ds4_group_mean_genero_p1_claridad_contenido.png}
\caption{ds4\_group\_mean\_genero\_p1\_claridad\_contenido.png}
\end{figure}
\begin{figure}[H]
\centering
\includegraphics[width=0.9\textwidth]{C:/sadi/static/plots/ds4_group_mean_genero_p2_dominio_docente.png}
\caption{ds4\_group\_mean\_genero\_p2\_dominio\_docente.png}
\end{figure}
\begin{figure}[H]
\centering
\includegraphics[width=0.9\textwidth]{C:/sadi/static/plots/ds4_group_mean_genero_p3_material_apoyo.png}
\caption{ds4\_group\_mean\_genero\_p3\_material\_apoyo.png}
\end{figure}
\begin{figure}[H]
\centering
\includegraphics[width=0.9\textwidth]{C:/sadi/static/plots/ds4_group_mean_genero_p4_participacion_clase.png}
\caption{ds4\_group\_mean\_genero\_p4\_participacion\_clase.png}
\end{figure}
\begin{figure}[H]
\centering
\includegraphics[width=0.9\textwidth]{C:/sadi/static/plots/ds4_group_mean_genero_p5_satisfaccion_general.png}
\caption{ds4\_group\_mean\_genero\_p5\_satisfaccion\_general.png}
\end{figure}
\section{Conclusión}
En conclusión, el análisis permite identificar una caracterización general del comportamiento de la muestra. La ausencia de asociaciones significativas destacadas sugiere cautela en la formulación de relaciones entre variables, siendo recomendable complementar estos hallazgos con análisis adicionales o con una revisión del diseño del instrumento. Asimismo, la identificación de variables clave permite focalizar la discusión en aquellos componentes que presentan mayor peso dentro del comportamiento observado en los datos. Este informe constituye una base preliminar de apoyo para la redacción de resultados, discusión académica y elaboración de informes técnicos o científicos.
\end{document}